\documentclass[11pt]{article}
\usepackage[a4paper,margin=27mm]{geometry}
\usepackage{amsmath,amssymb,amsthm,mathtools}
\usepackage{booktabs,graphicx,array,hyperref}
\hypersetup{colorlinks=true,linkcolor=blue,urlcolor=blue,citecolor=blue}
\newtheorem{theorem}{Theorem}[section]
\newtheorem{proposition}[theorem]{Proposition}
\newtheorem{lemma}[theorem]{Lemma}
\newtheorem{corollary}[theorem]{Corollary}
\newcommand{\E}{\mathrm E}
\newcommand{\N}{\mathbb N}
\title{Sharp Coordinate-Wise Plateaux in Six-Matching Covers\\
\large A bipartite construction, balanced vertex expansions, and certified lifting}
\author{Guillaume Sidoine\\Independent researcher\\
\href{https://orcid.org/0009-0005-4525-1182}{ORCID 0009-0005-4525-1182}}
\date{Research draft, 8 October 2026 --- not submitted or externally refereed}
\begin{document}
\maketitle
\begin{abstract}
For six perfect matchings of a cubic graph, let $Q$ be the sum of squared
edge-multiplicity deviations from two. We construct an explicit 18-vertex
3-edge-connected cubic bipartite graph with $Q=6$ at which no replacement
of one perfect matching strictly reduces $Q$. This is sharp: every
nonzero one-coordinate local minimum in a simple cubic graph has $Q\ge6$.
Moreover any state of a cubic bipartite graph with $0<Q\le10$
admits strict two-coordinate descent (three coordinates suffice for
$0<Q\le38$). At $Q=6$, two coordinates suffice to reach an exact twice-cover.
A balanced $K_{3,2}$ vertex expansion transports the obstruction to
3-edge-connected bipartite graphs of every order $18+4k$, maintaining
$Q=6$ and a two-coordinate escape. Bipartite matching dual certificates
also imply that local-optimum obstructions lift to arbitrary finite graph
covers, including connected covers of arbitrarily large girth (with the
potential growing with the covering degree). We distinguish these
coordinate-reconfiguration results from elementary consequences of
K\H onig's edge-colouring theorem and give a small, independently
checkable integer certificate. No unresolved covering conjecture is proved.
\end{abstract}

\section{Set-up and elementary barriers}
Let $G$ be a finite simple cubic graph with $n$ vertices, and let
$\mathcal M=(M_1,\ldots,M_6)$ be an ordered six-tuple of perfect matchings.
Put $m(e)=|\{i:e\in M_i\}|$ and
\[
 Q_G(\mathcal M)=\sum_{e\in E(G)}(m(e)-2)^2.
\]
Since each perfect matching covers $n/2$ edges and $G$ has $3n/2$
edges, $\sum_e(m(e)-2)=0$. At each vertex, the three incident
multiplicities sum to six. $Q=0$ means that every edge occurs exactly
twice. A \emph{one-coordinate local minimum} means that replacing any
one $M_i$ by any perfect matching cannot strictly reduce $Q$.
A \emph{$k$-coordinate correction} changes at most $k$ entries of the
ordered tuple, with no restriction on the distance between the old
and new perfect matchings.

\begin{lemma}\label{lem:lower}
For every simple cubic graph, if $0<Q\le4$, some one-coordinate
replacement strictly decreases $Q$. Hence every positive one-coordinate
local minimum has $Q\ge6$.
\end{lemma}
\begin{proof}
Write $d(e)=m(e)-2$. At each vertex the incident values of $d$ sum
to zero. Modulo two, $Q\equiv\sum_e d(e)=0$, so $Q$ is even.
If a nonzero edge has $|d(e)|\ge2$, its two endpoints each force
at least two additional units of squared defect among their other
incident edges. Simplicity ensures these are four different edges;
thus $Q\ge 4+4=8$. Consequently for $0<Q\le4$ every nonzero
defect is $\pm1$, and at each affected vertex one $+1$ and one
$-1$ edge meet. The defect edges therefore form disjoint cycles
with alternating signs. Each cycle has length at least four;
thus $Q=4$ and its support is an alternating 4-cycle.
The two $+1$ edges each occur in three matchings. Let their
sets of matching indices be $A$ and $B$, both of size three.
A $-1$ edge joining endpoints of these two positive edges
occurs in a matching whose index lies outside $A\cup B$;
therefore $A\cap B\ne\varnothing$. Choose $i\in A\cap B$.
On the 4-cycle, $M_i$ contains both positive edges and neither
negative edge. Flipping this alternating cycle in $M_i$ produces
a perfect matching and reduces $Q$ to zero.
\end{proof}

\section{An exact two-coordinate frontier}
\begin{proposition}[Frozen-factor completion]\label{prop:frozen}
Let $G$ be $d$-regular and bipartite, and let $t,d\ge1$ be integers.
Take $dt$ perfect matchings and fix a subcollection $F$ of them.
The fixed matchings can be completed to a $dt$-tuple in which every
edge occurs exactly $t$ times if and only if no edge occurs more
than $t$ times in $F$.
\end{proposition}
\begin{proof}
Necessity is immediate. For sufficiency put
$r(e)=t-\sum_{M\in F}\mathbf1_M(e)\ge0$.
At each vertex, $\sum_{e\ni v}r(e)=dt-|F|$, so the multigraph
formed by $r(e)$ parallel copies of edge $e$ is regular and bipartite.
K\H onig's edge-colouring theorem decomposes it into $dt-|F|$
perfect matchings, as required.
\end{proof}

\begin{lemma}[Small transversals for six matching labels]\label{lem:hit}
Let $S_1,\ldots,S_r$ be subsets of a six-element set, each of
size at least three. If $r\le5$, a two-element set meets every
$S_j$. If $r\le19$, a three-element set meets every $S_j$.
\end{lemma}
\begin{proof}
Every $S_j$ excludes at most three two-element subsets, since
its complement has size at most three. There are $15$ pairs in
total. If $r\le4$, at least one pair is not excluded. If $r=5$
and all pairs were excluded, equality in the counting bound
would force the complements of the five $S_j$ to be five
three-element sets partitioning the edge set of $K_6$ into
triangles. This is impossible because each vertex of $K_6$
has odd degree five, whereas any edge-disjoint union of
triangles has even degree at each vertex. Thus a pair exists
also when $r=5$. Likewise, each $S_j$ excludes at most
one three-element subset. Since there are $20$ triples, if
$r\le19$ at least one triple meets them all.
\end{proof}

\begin{theorem}[Bounded low-defect coordinated descent]\label{thm:range}
Let $G$ be a cubic bipartite graph and $\mathcal M$ a
six-matching state with $Q>0$. If $Q\le10$, a simultaneous
replacement of at most two matchings strictly decreases $Q$.
If $Q\le38$, at most three matchings suffice. In either case
such a replacement can be found by polynomial-time bipartite
matching computations (with at most six matching labels).
\end{theorem}
\begin{proof}
Put $d(e)=m(e)-2$. Because $\sum_e d(e)=0$,
\[
 Q=\sum_ed(e)^2\ge\sum_e|d(e)|
   =2\sum_{e:d(e)>0}d(e)
   \ge2|\{e:d(e)>0\}|.
\]
Orient surplus edges $d(e)>0$ from the left to the right
bipartition and deficit edges $d(e)<0$ from right to left.
At every vertex incident with a defective edge, the equation
$\sum_{e\ni v}d(e)=0$ guarantees at least one incoming
and one outgoing directed edge. Thus the defect digraph
contains a directed cycle $C$ with alternating surplus and
deficit edges. There are at most $Q/2$ surplus edges in the
entire graph. Each is present in at least three of the six
matching coordinates. When $Q\le10$ (respectively $Q\le38$),
Lemma~\ref{lem:hit} provides two (respectively three) matching
indices whose union contains every surplus edge on $C$.
Let $k$ be the number selected and let $x(e)$ be their current
aggregate multiplicity. Change $x(e)$ to $x(e)-1$ on surplus
edges of $C$ and to $x(e)+1$ on deficit edges of $C$, leaving
all other edges unchanged. This is nonnegative by the hitting
condition and preserves degree $k$ at every vertex, since
$C$ alternates. The modified integer multiplicity vector is
therefore a $k$-regular bipartite multigraph and decomposes
into $k$ perfect matchings by K\H onig. Replacing only the
selected matchings realizes this modification. For every
surplus or deficit edge on $C$, the absolute value of the
original defect drops by one; thus $Q$ strictly decreases.
Computationally, the hitting set can be found by testing the
fifteen pairs or twenty triples of six labels, and the factor
decomposition uses $k$ bipartite matching computations.
\end{proof}

\begin{theorem}[Sharp $Q=6$ coordination threshold]\label{thm:two}
For every simple cubic bipartite graph $G$, every six-matching
state with $Q=6$ admits a $2$-coordinate correction reaching $Q=0$.
There exists a 3-edge-connected cubic bipartite graph and a $Q=6$
state for which no $1$-coordinate correction reduces $Q$.
\end{theorem}
\begin{proof}
The argument in Lemma~\ref{lem:lower} shows that $Q=6$ cannot
contain an edge with $|d(e)|\ge2$. Hence the support of $d$ consists
of one alternating 6-cycle. Its three positive edges each occur
in exactly three of the six matchings. By Lemma~\ref{lem:hit}, choose
two matching indices that together meet all three positive edges.
Freeze the other four matchings. Every edge outside the cycle has
total multiplicity two, every negative edge has total multiplicity
one, and on each positive edge at least one copy lies in the two
unfrozen matchings. Thus no edge is used more than twice among
the frozen four matchings. Proposition~\ref{prop:frozen} with
$(d,t)=(3,2)$ completes the frozen four to an exact twice-cover
by changing only the selected pair. The second assertion is proved
by the explicitly certified example below.
\end{proof}

\begin{proposition}[Nonbipartite $Q=6$ pair criterion]
For a six-matching state on any simple cubic graph with $Q=6$,
a direct two-coordinate correction to $Q=0$ exists if and only
if there are two matching indices that together meet all three
surplus edges, and after freezing the other four matchings the
residual multiplicity vector decomposes into even cycles
(counting doubled edges as cycles of length two).
\end{proposition}
\begin{proof}
For a chosen pair $I$, let $f(e)$ count its uses among the
other four (frozen) matchings, and put $r(e)=2-f(e)$.
The vector $r$ is nonnegative exactly when $I$ meets every
surplus edge; this uses the alternating-six-cycle structure
at $Q=6$. Independently, $\sum_{e\ni v}r(e)=2$ at every vertex.
Thus the residual is a spanning degree-two multigraph.
It is the union of two perfect matchings if and only if all
its cycle components have even length. This equivalence does
not require bipartiteness of the ambient cubic graph.
\end{proof}

\begin{proposition}[Exact two-coordinate block optimization]
For any fixed pair of indices in a sextuple of perfect matchings
of a bipartite cubic graph, the best simultaneous replacement
of those two matchings minimizing $Q$ can be found in
polynomial time by an integer minimum-cost flow.
\end{proposition}
\begin{proof}
Freeze four matchings and denote their edge multiplicities by
$f(e)$. Choose a new residual multiplicity $x(e)\in\{0,1,2\}$
such that $\sum_{e\ni v}x(e)=2$ at every vertex, minimizing
$\sum_e(f(e)+x(e)-2)^2$. For each edge, the incremental costs
of the first and second residual copies are
$2(f(e)-2)+1$ and $2(f(e)-2)+3$, respectively. Encode these
as two unit-capacity parallel marginal arcs in a bipartite
transportation network requiring throughput two at each
vertex. Its integer minimum-cost flow is optimal; the increasing
marginal costs ensure the second copy is never chosen without
the first in an optimal flow. The resulting degree-two bipartite
multigraph splits into two perfect matchings by K\H onig's
theorem, giving an optimal pair of replacement coordinates.
\end{proof}

\section{The 18-vertex certificate}\label{sec:certificate}
The bipartition is $L=\{0,\ldots,8\}$ and
$R=\{9,\ldots,17\}$. Number the following edges in the order
shown, starting from index $0$:
\begin{center}\small
\begin{tabular}{lll}
$(0,11),(0,14),(0,9)$ & $(1,14),(1,13),(1,11)$ & $(2,15),(2,17),(2,13)$\\
$(3,16),(3,11),(3,17)$ & $(4,9),(4,15),(4,12)$ & $(5,13),(5,12),(5,10)$\\
$(6,17),(6,10),(6,15)$ & $(7,12),(7,16),(7,14)$ & $(8,10),(8,9),(8,16)$
\end{tabular}
\end{center}
Represent matchings as $27$-bit nonnegative integers using
these edge indices. The initial ordered six-tuple is
\[
(76093521,26264228,19173986,70392073,38347922,21308684).
\]
Its six defective edges are, with multiplicities,
\[
\begin{array}{c|rrrrrr}
e&(4,9)&(4,12)&(5,12)&(5,10)&(8,10)&(8,9)\\ \hline
m(e)&3&1&3&1&3&1
\end{array}
\]
so $Q=6$. The path obtained by replacing the third matching
by $36053794$, followed by replacing the sixth matching
by $21283916$, has potentials $6\to6\to0$.

For an independent universal certificate of local optimality,
fix $i$ and put
\[
 w_i(e)=2(m(e)-2)+1-2\mathbf1_{M_i}(e).
\]
For any other perfect matching $N$ one checks exactly that
\[
 Q(M_1,\ldots,N,\ldots,M_6)-Q(\mathcal M)
 = \sum_{e\in N}w_i(e)-\sum_{e\in M_i}w_i(e).\tag{$\ast$}
\]
The supplementary JSON file supplies, for each $i$, integers
$a_i(v)$ at all eighteen vertices such that
\[
 a_i(u)+a_i(v)\le w_i(uv)\quad (uv\in E),\qquad
 a_i(u)+a_i(v)=w_i(uv)\quad(uv\in M_i).
\]
Summing the inequalities over any perfect matching $N$ and
using the equalities for $M_i$ proves that $(\ast)\ge0$.
These six dual witnesses are each 18 integers; the
\texttt{standalone\_checker.py} file checks all inequalities using
only Python's standard library. It also verifies that the graph
stays connected upon deleting any zero, one or two edges,
certifying 3-edge-connectivity. A second implementation
exhaustively enumerates all $47$ perfect matchings and finds
minimum replacement change zero in every coordinate.

\section{Lifting a sharp plateau to an infinite family}
\begin{lemma}[Balanced vertex expansion]\label{lem:expand}
Let $G$ be a simple cubic graph with a six-matching state
$\mathcal M$ and let $v$ have incident multiplicities $(2,2,2)$.
Replace $v$ with the 3-pole $K_{3,2}$, attaching the three
former neighbours to its three degree-two vertices.
The resulting cubic graph $G'$ has a state $\mathcal M'$ with
$Q_{G'}(\mathcal M')=Q_G(\mathcal M)$. If $\mathcal M$ is a
one-coordinate local minimum, so is $\mathcal M'$.
This operation preserves simplicity, bipartiteness when present,
and 3-edge-connectivity when present.
\end{lemma}
\begin{proof}
Each of the three incident edges at $v$ occurs in exactly two
of the six matchings. Every matching chooses one incident edge.
In $K_{3,2}$, the two ports not chosen externally must be
matched to the two internal degree-three vertices. For the two
matching labels choosing the same external port, select opposite
internal assignments. Each of the six internal edges then occurs
exactly twice; every other edge retains its old multiplicity.

Every perfect matching of $G'$ restricts to a perfect matching
of $G$ by contracting the 3-pole: both internal hubs must be
matched to ports, leaving exactly one external port matched.
Replacing one lifted matching changes the original-edge
contribution to $Q$ by an amount at least zero whenever the
original tuple is a local minimum. The initial internal
contribution is zero, and every new internal contribution is
nonnegative. This proves the claim.

The gadget adds four vertices and is bipartite with its three
ports in the old vertex's part. To check 3-edge-connectivity,
consider any cut of size at most two. If all three ports are
on the same side but an internal hub is on the other side,
at least three internal edges cross. If the ports are split,
at least two internal edges cross; equality two forces both
hubs with the majority ports. A cut of size at most two then
has no external crossing edges, so contracting the gadget
would yield a cut of size one in $G$ (the minority port's
external neighbour lies on the minority side), a contradiction.
If the whole gadget lies on one side, the cut corresponds
unchanged to a cut in $G$. Hence no such cut exists.
\end{proof}

\begin{corollary}[Arbitrarily large, 3-edge-connected sharp traps]
For every $k\ge0$ there is a connected 3-edge-connected cubic
bipartite graph on $18+4k$ vertices with six perfect matchings
forming a positive one-coordinate local minimum of defect
exactly $Q=6$. The tuple admits a $2$-coordinate correction
to $Q=0$; indeed the base $Q=6\to6\to0$ sequence lifts
unchanged along the family.
\end{corollary}
\begin{proof}
The vertex $0$ in the base graph has incident multiplicities
$(2,2,2)$, and the two replacements in Section~\ref{sec:certificate}
do not change its matching choices. Apply
Lemma~\ref{lem:expand}; the internal hubs are again balanced
and their incident matching choices stay fixed during both
replacements, so the construction can be repeated indefinitely.
The local minimum and defect persist, as does the indicated path.
\end{proof}

\section{Lifting under finite covers}
\begin{proposition}[Covering stability of bipartite dual certificates]
Let $G$ be cubic and bipartite, and let $\mathcal M$ be a
one-coordinate local minimum for $Q$. Every finite
$k$-sheeted graph covering $p:H\to G$ inherits a one-coordinate
local minimum $p^{-1}(\mathcal M)$ with potential $kQ_G(\mathcal M)$.
\end{proposition}
\begin{proof}
The preimage of a perfect matching is a perfect matching, and
edge multiplicities pull back along $p$, so the potential scales
by $k$. For each coordinate $i$, the minimum-weight perfect
matching problem with weights $w_i$ has an integral bipartite
assignment dual attaining the original matching's weight.
Pull its vertex potentials back along $p$. The edge inequalities
and the equalities on the lifted matching remain valid because
every edge of $H$ projects to an edge of $G$. Thus no lifted
coordinate can strictly improve the potential.
\end{proof}

\begin{corollary}[High-girth obstructions]
There are connected bridgeless cubic bipartite graphs of
arbitrarily large girth that possess positive one-coordinate
local minima of the six-matching quadratic defect satisfying
$Q=|V(G)|/3$.
\end{corollary}
\begin{proof}
A finite connected graph has finite connected covers of
arbitrarily large girth: the fundamental group is finitely
generated free, hence residually finite; taking a finite-index
normal subgroup excluding the finitely many nontrivial closed
reduced walks of length at most a chosen bound produces such
a cover. Apply the preceding proposition to the base certificate.
Every edge of a finite cover of a bridgeless graph lies on a
cycle (a lift of some positive multiple of a cycle downstairs),
so the covers remain bridgeless. The base graph has 18 vertices
and defect six; both quantities scale with the covering degree.
\end{proof}

An explicit connected $2$-fold cover obtained by twisting
edge indices $0$ and $6$ has $36$ vertices, girth $6$, edge
connectivity $3$, and a locally optimal tuple with $Q=12$.
The provided six lifted dual certificates check this example.

\section{Scope, relation to prior work, and outlook}
The frozen-factor completion criterion is an elementary
application of K\H onig's theorem and is not presented as a
new existence theorem for cubic bipartite graphs. Similarly,
our coordinate changes may replace entire perfect matchings:
they are not 4-cycle flips in the sense of standard perfect
matching reconfiguration \cite{Bonamy2019}. A recent treatment
of toric ideals and $3$-colour transformations for \emph{fixed}
edge multiplicities \cite{Mori2026} is related to neutral-fibre
navigation, whereas the present obstruction concerns descent of
an explicit quadratic multiplicity defect. Recent approaches to
the Berge--Fulkerson conjecture through partial coverings are
also distinct \cite{Luo2026}.

The main contribution asserted here is the exact sharp
one-versus-two coordinate obstruction at the first positive
defect level together with the infinite lifting results.
Whether this qualifies for a research note requires a targeted
literature audit and independent external specialist review.
Nothing here proves Berge--Fulkerson for nonbipartite cubic graphs.
Promising follow-ups include determining the minimum coordinate
width for strict descent in nonbipartite graph classes, and
using small dual certificates as Lean-checkable obstruction witnesses.

\paragraph{Reproducibility.} All graph edges, ordered matching masks,
six integer dual potentials, the stand-alone verifier, the
$K_{3,2}$ lift implementation, the covering construction, and
regression tests accompany this draft. The current environment
contains no Lean 4 installation; formalization has therefore
not been completed or claimed.

\paragraph{Research-process disclosure.} Computer-assisted search,
certificate generation, proof exploration and drafting used
ChatGPT (GPT-6); an independently checkable proof record is
provided, and the draft remains subject to external review.

\begin{thebibliography}{9}
\bibitem{Konig} D. K\H onig, \emph{Gr\'afok \text{\'es} alkalmaz\'asuk a determin\'ansok \text{\'es} a halmazok elm\'elet\'ere},
Matematikai \text{\'es} Term\'eszettudom\'anyi \text{\'E}rtes\'it\H o \textbf{34} (1916), 104--119.
\bibitem{Bonamy2019} M. Bonamy, N. Bousquet, M. Heinrich,
T. Ito, Y. Kobayashi, A. Mary, M. M\"uhlenthaler and K. Wasa,
\emph{The perfect matching reconfiguration problem},
MFCS 2019, LIPIcs \textbf{138}, Art. 80.
\href{https://doi.org/10.4230/LIPIcs.MFCS.2019.80}{doi:10.4230/LIPIcs.MFCS.2019.80}.
\bibitem{Mori2026} K. Mori, R. Motomura, H. Ohsugi and A. Tsuchiya,
\emph{Toric ideals of matching polytopes and edge colorings},
Mathematika (2026).
\href{https://doi.org/10.1112/mtk.70099}{doi:10.1112/mtk.70099}.
\bibitem{Luo2026} Y. Luo, R.-X. Hao, R. Luo, C.-Q. Zhang and W. Zhou,
\emph{Berge--Fulkerson conjecture, perfect matching partial coverings and odd dividers},
Journal of Graph Theory (2026).
\href{https://doi.org/10.1002/jgt.70097}{doi:10.1002/jgt.70097}.
\end{thebibliography}
\end{document}
